\chapter{Conclusions}
\label{ch:conclusions}

\section{Summary of the study and its results}
\label{sec:concl-answer}

The study investigated which core designs a LABGENE-class pressurised
water reactor of \SI{48}{MW} thermal could adopt when the only
information available is the open literature. It answered by building
a surrogate-assisted multi-objective optimisation around a Monte Carlo
model of the core and executing it over 9 campaigns, 636 evaluations,
each campaign changing one element of the method, followed by a
post-analysis of Campaign 9 on the three-dimensional core and a
continuation of its search with 36 further evaluations. Campaigns 1 to
8 maximised the cycle length and minimised the radial peaking factor.
Campaign 9 minimised the peaking factor and the critical boron concentration under a minimum cycle length. Within the model and the design space of this work, the study reached five results:

\begin{enumerate}
  \item \textbf{In this study a limit derived from a few designs was not
        a substitute for the measurement on each design:} in the four
        cases where one replaced the other, the derived limit
        misclassified designs near the boundary it decided:
        \begin{itemize}
          \item the upper reactivity bound of 1.126 on the core
                multiplication factor, derived from the bank worth of 8
                Campaign 6 designs, rejected 14 designs whose own bank
                worth made them controllable
                (Section~\ref{sec:res-c7-screens}),
          \item the MTC boron limit of C8-47, applied to every
                Campaign 9 design, accepts C9-57, which its own limit
                rejects by \SI{193}{ppm} (Section~\ref{sec:res-c9-ceiling}),
                    \item the constant minimum assembly cycle length of stage 1
                of the continuation admitted C9-72, which fails the
                mission by \SI{44}{d} (Section~\ref{sec:res-c9-cont}),
          \item the gadolinia-dependent minimum assembly cycle length of
                stage 2 admitted C9-85 and C9-88, which fail, and
                rejected C9-11, which passes
                (Section~\ref{sec:res-c9-cont}).
        \end{itemize}
  \item \textbf{The cycle length of this core is a requirement to be
        satisfied and not an objective:} every Campaign 8 front member
        satisfied the refuelling interval of \SI{1826}{EFPD}, with
        margins from \SI{115}{EFPD} to a factor of 3.2, and the excess reactivity of
        the longer cycle, \num{8930} to \SI{13627}{pcm}, had to be
        compensated by the four regulating banks at a large insertion
        depth at power (Section~\ref{sec:concl-reform}), so the
        challenge of this core is the control of the excess reactivity
        that a longer cycle requires.
  \item \textbf{At the required cycle length the boron the core needs
        decreases as its gadolinia loading increases:} among the
        feasible Campaign 9 designs within \SI{0.35}{wt\%} of
        \SI{4.31}{wt\%} enrichment, the critical boron concentration at
        the operating maximum decreases from 3003 to \SI{1406}{ppm} as
        the gadolinia increases from 0.14 to \SI{4.42}{wt\%}, because a
        concentrated pin burns out slowly enough that no mid-cycle hump
        appears (Section~\ref{sec:res-c9-gd}).
    \item \textbf{The assembly depletion, with its axially uniform
        burnup, overestimates the cycle length of the core:} the
        three-dimensional core depleted in 8 axial layers has a cycle 15
        to 25 per cent shorter than the assembly value, while the same
        core depleted with uniform burnup is within a few per cent of
        it, so the overestimate comes from the burnup distribution along
        the height (Section~\ref{sec:res-c9-axial}). Its effect depends
        on the formulation:
        \begin{itemize}
          \item with the cycle length as an objective, the overestimate
                shifts the front towards longer cycles and, because the
                reduction depends on the gadolinia loading, can reorder
                it,
          \item with the cycle length as a constraint, the front is
                placed on that constraint, so the reduction decides which
                designs are feasible and has to be applied to the
                constraint. No member of the Campaign 9 front satisfies
                the minimum cycle length once it is applied, and near
                that minimum only the eight-layer depletion applies it
                with enough accuracy.
        \end{itemize}
  \item \textbf{A proxy has to be validated on the ordering it produces near the front:} the single-seed estimator of Campaign 2, the
        assembly peaking factor of Campaign 3 and the two-dimensional
        peaking factor of Campaign 9 pass the correlation usually
        applied to them and still reorder the designs near the front
        (Sections~\ref{sec:res-corevalidation}
        and~\ref{sec:res-c9-peaking}).
\end{enumerate}

\subsection{The proxies of the study and their validation}
\label{sec:concl-proxies}

Each campaign evaluated its designs through proxies, quantities computed in place of the ones wanted, because the wanted ones were too expensive to
compute at every evaluation. The study used 13 proxies. Ten
are used inside the optimisation loop, where they decide which designs
are evaluated and which are feasible, and the other 3, the linear boron
extrapolation, the lattice moderator temperature coefficient and the
common MTC boron limit, enter the post-analyses. Twelve were tested and
are listed below, and the thirteenth, the depletion at a constant
\SI{1000}{ppm}, was never tested.

Each proxy was compared with the quantity it approximates, computed by
the more expensive calculation it replaces, for example the assembly
peaking factor with the core solve and the assembly cycle length with
the eight-layer depletion, and the comparison decided whether it was
kept or replaced. Five proxies remained within a stated magnitude:

\begin{itemize}
  \item \textbf{Leakage-corrected $k_\mathrm{target}$ on the assembly
        depletion, every campaign, for the end of cycle of the core:}
        on 10 Campaign 8 designs the drift over burnup is
        $-13 \pm \SI{30}{pcm}$, and the beginning-of-life residual,
        $+174$ to \SI{-329}{pcm} ordered by enrichment, is worth $-14$
        to \SI{+56}{d} of cycle without changing the order of the front
        (Section~\ref{sec:res-kt-burnup}). Kept, since a core depletion
        at every evaluation is not affordable.
  \item \textbf{Axial leakage factor $L_\mathrm{ax} = 1.0289$,
        Campaigns 8 and 9, for the unrodded finite-height core:} within
        $-25$ to \SI{+211}{pcm} on the Campaign 8 designs and 80 to
        \SI{240}{pcm} above it on the Campaign 9 designs measured in
        three dimensions, against a
        systematic \SI{80}{pcm}. The rodded factors are smaller, so the
        two-dimensional controllability constraint is conservative by
        1272 to \SI{2671}{pcm} (Section~\ref{sec:res-kt-burnup}). Kept
        for the unrodded multiplication factor.
  \item \textbf{Gaussian-process surrogates, every campaign, for the
        evaluator inside the search:} leave-one-out $R^2$ of 0.988 for
        the cycle length and 0.690 for the peaking factor in
        Campaign 8, and no predictive accuracy for the cycle length on
        the feasible designs of Campaign 9, $R^2 = -0.41$, where that
        quantity is an active constraint
        (Sections~\ref{sec:res-c8-validation}
        and~\ref{sec:res-c9-validation}). Kept, because the surrogates
        preserve the ordering of the design space, with $R^2 = 0.883$
        for the cycle length over all 60 Campaign 9 designs, so the
        search on them moves towards the region of better designs. 
  \item \textbf{Linear extrapolation of the critical boron
        concentration from the solve at \SI{1000}{ppm}, Campaign 8
        post-analysis, for the root of the reactivity in boron:} within
        \SI{21}{ppm} of the root on C8-47, the one design it was
        checked on (Section~\ref{sec:res-c8-boron}). Replaced in
        Campaign 9 by the root found between 2 solves for every design.
  \item \textbf{Lattice moderator temperature coefficient, Campaign 9,
        for the coefficient of the core:} measured on 11 lattices, with
        the gadolinia loading confounded with the reflector thickness,
        so the fitted limit is valid inside the converged region of the
        campaign only (Section~\ref{sec:res-c9-ceiling}). Kept as the
        MTC boron limit of each lattice, outside the loop.
\end{itemize}

Seven proxies failed on the ordering or the bound they were meant to
preserve:

\begin{itemize}
  \item \textbf{Single-seed assembly estimator at $4\,000\times60$,
        Campaigns 1 and 2, for the converged assembly peaking factor:}
        stored values on average 0.07 below the high-fidelity
        reference against a seed-to-seed standard deviation of
        $0.015 \pm 0.003$ and a rank inversion of the finalists at
        $|t| = 12.8$ (Sections~\ref{sec:res-c1c2}
        and~\ref{sec:res-c3-fidelity}). Replaced from Campaign 3 by a
        seed per design and $16\,000\times120$.
  \item \textbf{Reflected-assembly peaking factor, Campaigns 1 to 3,
        for the core peaking factor:} Spearman coefficient of $+0.890$
        over the 36 uniform cores of Campaign 3 and a reordering of up
        to 10 places near the front
        (Section~\ref{sec:res-corevalidation}). Replaced by the core
        solve as the objective from Campaign 4.
  \item \textbf{Derived upper reactivity bound, Campaign 7, for
        controllability by the regulating banks:} the bound of 1.126
        rejected 14 designs whose measured bank worth made them
        controllable (Section~\ref{sec:res-c7-screens}). Kept as an
        outer bound from Campaign 8, with the measured controllability
        constraint binding.
  \item \textbf{Two-dimensional radial peaking factor, Campaigns 4 to
        9, for the peaking factor of the three-dimensional core:} the
        three-dimensional value shifts the 5 Campaign 9 front designs
        by $+0.006$ to $-0.030$, and the shift changes which designs
        are non-dominated, C9-40 entering the front and C9-44 and C9-35
        leaving it (Section~\ref{sec:res-c9-peaking}). Kept as the
        objective of the search, with the peaking of every front
        measured in three dimensions in the post-analysis.
  \item \textbf{Reflective-assembly depletion with axially uniform
        burnup, every campaign, for the cycle length of the finite
        core:} the eight-layer depletion of the three-dimensional core
        gives a cycle 15 to 25 per cent shorter than the
        reflective-assembly depletion on the 18 designs depleted both
        ways. C9-47 and C9-35, the two Campaign 9 front members that
        satisfy the minimum cycle length with the reflective-assembly
        depletion, are 253 and \SI{185}{d} short of it with the
        eight-layer depletion (Section~\ref{sec:res-c9-axial}). Replaced, for every front
        member, by the eight-layer depletion of the three-dimensional
        core, and the front rescored on it and the search continued
        under it give C9-27, C9-69 and C9-70.
  \item \textbf{MTC boron limit of C8-47, \SI{2763}{ppm}, Campaign 9,
        for the limit of each lattice:} lattice limits of 2091 to
        \SI{3128}{ppm} on 11 lattices, and the common limit accepts
        C9-57, which its own limit rejects by \SI{193}{ppm}
        (Section~\ref{sec:res-c9-ceiling}). Replaced by the limit of
        each candidate lattice, and kept only for designs whose own
        limit was not calculated.
    \item \textbf{Gadolinia-dependent minimum assembly cycle length,
        Equation~\eqref{eq:axial-floor}, stage 2 of the continuation of
        Campaign 9, for the eight-layer depletion:} leave-one-out error
        of 38 to \SI{51}{d} against margins of 7 to \SI{13}{d} near the
        mission cycle length of \SI{1826}{EFPD}. The error of the rule is 3 to 6 times the margins by
        which the designs near the minimum cycle length pass or fail,
        so it cannot decide their feasibility, and it admits C9-85 and
        C9-88, which fail when measured, and rejects C9-11, which
        passes (Section~\ref{sec:res-c9-cont}).
\end{itemize}

The untested proxy, the depletion at a constant \SI{1000}{ppm}, is the proxy for the boron letdown of the operated core,
so every cycle length of the study is a lower bound. The tests apply
theory that the earlier chapters cite: the correlated sampling of Lux
and Koblinger~\cite{lux_koblinger_1991} for the seed effect, Spearman's
coefficient~\cite{kendall_gibbons_1990} and the multiplicative
discrepancy of Kennedy and O'Hagan~\cite{kennedy_ohagan_2000} for the
ordering and bridge tests, the multiplication-factor description of
Duderstadt and Hamilton~\cite{duderstadt_hamilton_1976} for the leakage
hierarchy, and the leave-one-out and walk-forward procedures of
Section~\ref{sec:validation-framework} on the Gaussian process of
Rasmussen and Williams~\cite{rasmussen_gp_2006}. The derived bound, the
lattice coefficient and the constant boron have no basis in the
literature cited and are assumptions of this work.

\section{The two formulations of the optimisation problem}
\label{sec:concl-reform}

Campaigns 1 to 8 maximised the cycle length against the peaking factor,
a reasonable formulation for a core whose attainable cycle length is
not known in advance. Campaign 8 showed its restriction once
controllability was measured. Every front member satisfied the
refuelling interval of \SI{1826}{EFPD} with margin, the objective kept increasing
the cycle length after the requirement was satisfied, and it could do
so only by raising the beginning-of-life excess reactivity, \num{8930}
to \SI{13627}{pcm}, which the four regulating banks then compensate at
a large insertion depth at power (Section~\ref{sec:res-c8-mission}).
The two-bank designs of the same campaign, at \num{2564} to
\SI{4068}{pcm}, are in the regime of the NuScale-like reference core,
whose rods compensate about \SI{2700}{pcm}, and 2 of them, C8-13 and
C8-31, satisfy the interval (Section~\ref{sec:res-c8-split}).

Campaign 9 therefore minimised the peaking factor and the critical
boron concentration at the operating maximum, with the minimum cycle
length as a constraint (Section~\ref{sec:meth-c9}). The second
objective expresses the beginning-of-life excess reactivity and the
gadolinium hump in one quantity, in the unit in which the moderator
temperature coefficient limits it. The reformulation follows from three
results of Campaign 8:

\begin{itemize}
  \item \textbf{Cycle length:} in a design space with enrichment up to
        \SI{13.9}{wt\%}, every front member reaches the minimum cycle
        length of the mission with margin, so the cycle length no
        longer discriminates between candidates.
  \item \textbf{Excess reactivity:} the longer cycle is obtained only by
        a larger excess reactivity at beginning of life, which is the
        quantity that limits the design.
  \item \textbf{Control rods:} the bank design and the rod worths of
        LABGENE are not public, so the rod worths and the MTC boron
        limit are those of the NuScale-like reference. Within them, an
        objective on the critical boron concentration selects the
        designs whose excess reactivity boron and burnable absorber
        compensate, with only the first two regulating banks in normal
        operation. None of the 8 Campaign 8 front members satisfies
        that criterion, while every member of the Campaign 9 front and
        of the final front does (Section~\ref{sec:res-c9-ctrl3d}).
\end{itemize}

The results of Campaigns 8 and 9 determine the domain of application
of each formulation:

\begin{itemize}
  \item \textbf{Cycle length as an objective, Campaigns 1 to 8:} the
        better formulation when the attainable cycle length of a design
        space is unknown, since it measures how far the cycle can be
        extended and at what peaking. It has no bound on the excess
        reactivity, so its front is controllable only with the
        regulating banks deep in the core.
  \item \textbf{Cycle length as a constraint, Campaign 9:} the better
        formulation once the mission fixes the cycle length, since a
        cycle beyond the requirement gives the mission nothing and
        requires excess reactivity that the control system must
        compensate. Its front is at the minimum cycle length and is
        operable with boron and burnable absorber. The campaign front
        needs 1406 to \SI{1908}{ppm} at the operating maximum, and the
        first two regulating banks alone keep every member subcritical
        by 9228 to \SI{9545}{pcm} in three dimensions
        (Section~\ref{sec:res-c9-ctrl3d}). The final front C9-27, C9-69
        and C9-70, after the axial correction, needs 1633 to
        \SI{2268}{ppm}, with two-bank margins of 2222 to \SI{6316}{pcm}
        (Section~\ref{sec:res-c9-cont-front}).
\end{itemize}

\section{Achievement of the objectives of the study}
\label{sec:concl-candidates}

The first objective of Section~\ref{sec:intro-objectives}, the
framework, was reached. Its loop has two parts
(Section~\ref{sec:meth-loop}):

\begin{itemize}
  \item \textbf{Search on the surrogates:} the NSGA-II search on the
        Gaussian-process surrogates proposes the designs to evaluate,
        and it determines the direction of the search, the region of
        the design space where the surrogates predict the front to be.
        Figure~\ref{fig:c9-gp-fronts} validates this part in Campaign 9:
        as the training data grow from 24 to 60 evaluations, the
        predicted front narrows in $F_{\Delta H}$ from 1.487 to 1.640 to
        1.495 to 1.535, onto the 1.490 to 1.545 of the evaluated front,
        so the surrogates give the search its direction.
  \item \textbf{OpenMC evaluation:} the transport evaluation of the
        selected batch computes the objectives and constraints of those
        designs, and it determines the accuracy of the front, since
        every value of a front member is a transport result.
\end{itemize}

 The acquisition
function of the first part selects, among the predicted front, the
designs whose evaluation reduces the uncertainty of the surrogates
most, and their results enter the training data at the next iteration,
so with each infill iteration the surrogates are refitted on more
evaluations near the front and their predictions there become more
precise, until the stopping rule of Section~\ref{sec:meth-loop-stop}
finds no further gain in hypervolume.

The framework gave a Pareto front in every campaign, and the quality
of that front depends on the refinement of the model and on the
definition of the constraints, so each campaign from 1 to 9 improved
the optimiser by changing one element of it. Campaign 6 defined the
acquisition function one element per infill block, and Campaign 9
showed the restriction of that function when the search is on a
constraint that is active on the front, where its feasibility margin
rejects every candidate (Section~\ref{sec:res-c9-acq}). The surrogates
and the fronts of the last two campaigns were validated against the
evaluations and against the transport noise
(Sections~\ref{sec:res-c8-validation}
and~\ref{sec:res-c9-validation}).

The second objective of Section~\ref{sec:intro-objectives}, a set of
candidate core configurations, was also reached. It is the result of
the post-analysis of Campaign 9 on the three-dimensional core depleted
in 8 layers and of the continuation of the search under that cycle
length, the most refined optimiser and core model of this thesis, and
it replaces the candidates of the campaigns. The candidate
configurations are the front C9-27, C9-69 and C9-70
(Section~\ref{sec:res-c9-cont-front}), the most accurate results of
the study. That measurement is also the validation
against a model of higher fidelity than the loop model that the first
objective requires. Table~\ref{tab:concl-front} gives every
quantity measured on them.

\begin{table}[htbp]
  \centering
  \small
  \begin{threeparttable}
  \caption{Final front of Campaign 9: design variables and three-dimensional measurements.}
  \label{tab:concl-front}
  \setlength{\tabcolsep}{4pt}
  \begin{tabular}{@{}lccc@{}}
    \toprule
    Quantity & C9-27 & C9-69 & C9-70 \\
    \midrule
    Enrichment $e$ [wt\%] & 5.11 & 4.87 & 4.87 \\
    Gadolinia $w_\mathrm{Gd}$ [wt\%] & 5.70 & 4.38 & 4.41 \\
    Poisoned pins $n_\mathrm{Gd}$ & 24 & 16 & 20 \\
    Reflector thickness $t_\mathrm{refl}$ [cm] & 5.66 & 5.66 & 5.66 \\
    $F_{\Delta H}$, 2D campaign value & 1.603 & 1.526 & 1.557 \\
    $F_{\Delta H}$, 3D & 1.565 & 1.496 & 1.516 \\
    $c_\mathrm{max}$ [ppm] & 1633 & 2268 & 1840 \\
    Assembly hump $\Delta\rho_\mathrm{hump}$ [pcm] & $-1800$ & $-1579$ & $-1010$ \\
    $\mathrm{EFPD}_\mathrm{3D}$, 8 layers [d] & $1926 \pm 10$\tnote{a} & $1895 \pm 7$ & $1850 \pm 3$\tnote{a} \\
    Margin over \SI{1826}{d} [d] & $+100$ & $+69$ & $+24$ \\
    Axial leakage factor $L_\mathrm{ax}$ & 1.030 & 1.031 & 1.030 \\
    Two-bank margin, 3D, \SI{1000}{ppm} [pcm] & 6316 & 2222 & 4941 \\
        MTC boron limit at \SI{12.8}{MPa} [ppm] & $3128 \pm 111$\tnote{b} & $2867 \pm 115$ & $3018 \pm 102$\tnote{b} \\
    MTC margin [ppm] & $+1495$ & $+599$ & $+1178$ \\
    Origin & rescoring of Campaign 9 & stage 1 & stage 1 \\
    \bottomrule
  \end{tabular}
  \begin{tablenotes}[flushleft]\footnotesize
    \item[a] Mean of 5 seeds and standard error of the mean.
        \item[b] Extrapolated beyond the scanned \SI{3000}{ppm}.
  \end{tablenotes}
  \end{threeparttable}
\end{table}

C9-27 is the recommended design. It needs the least boron of any design
that satisfies the minimum cycle length, has the largest cycle margin,
the largest MTC margin and the largest two-bank margin of the front, and
no mid-cycle hump. C9-69 and C9-70 give a lower peaking
factor at 635 and \SI{207}{ppm} more boron and a smaller margin.

\section{Limitations}
\label{sec:concl-limits}

Section~\ref{sec:res-uncertainty} collects the sources of uncertainty
of the study, and 7 of them affect the conclusions:

\begin{enumerate}
  \item \textbf{Axially uniform burnup in the loop:} every campaign
        cycle length is an assembly value that overestimates the core
        by 15 to 25 per cent. The final front comes from the
        post-analysis and from a continuation that applied the
        eight-layer depletion to front members only, and stage 2 spent
        10 of its 18 evaluations above the MTC boron limit, so a better
        design under the corrected requirement may exist in the region
        the search did not reach.
  
  \item \textbf{Two-dimensional peaking in the loop:} the campaign
        value overestimates the peaking factor by 0.060 to 0.077, and
        the three-dimensional shift, up to $-0.030$, reorders
        neighbouring designs, so 5 of the 10 pairs of the campaign
        front are not distinguishable by one campaign solve.
  \item \textbf{Constant boron and no thermal-hydraulic feedback:}
        every cycle length is a lower bound for a core with boron
        letdown, and the moderator temperature coefficient enters only
        through a limit measured outside the loop on 11 lattices, valid
        inside the converged region of Campaign 9, with margins
        conservative by about \SI{450}{ppm} because the limit is
        three-dimensional and the concentration two-dimensional.
  \item \textbf{MTC boron limit outside the search:} the limit was
        stored and not constrained, so the acquisition placed 6 of the
        36 infill evaluations of Campaign 9 and 10 of the 18 of stage 2
        in the region below \SI{1}{wt\%} of gadolinia that the limit
        excludes.
    \item \textbf{Reference fuel loading of $k_\mathrm{target}$:} the
        radial leakage factor decreases by $-64 \pm \SI{9}{pcm}$ per
        wt\%, so the $k_\mathrm{target}$ read from the table, computed on
        the reference fuel loading, is too high for a design of higher
        enrichment, by about 500 to \SI{600}{pcm} at the upper
        enrichment bound of Campaign 8. The depletion of such a design
        reaches the end-of-cycle criterion earlier than its core would,
        so its reported cycle length is shorter than the true one, which
        is conservative. A new campaign that maximises the cycle length,
        as Campaigns 1 to 8 did, should therefore compute
        $k_\mathrm{target}$ over the enrichment as well as the reflector
        thickness.
    \item \textbf{Operating rather than licensing criteria:} two criteria
        of this work are operating criteria and not the ones a licensing
        analysis applies:
        \begin{itemize}
          \item the controllability constraint tests the rodded core at
                hot full power with the operating boron, and not the
                shutdown margin,
          \item the discharge-burnup limit acts on the assembly-average
                burnup, while fuel performance limits apply to the local
                burnup, which reaches 1.35 times the core average in the
                central layers of the axially resolved depletion of
                Section~\ref{sec:res-c9-axial} and is not constrained.
        \end{itemize}
  \item \textbf{Boundary optimum:} every front member has the largest
        reflector the vessel admits, so the front is conditional on the
        vessel envelope and the clearance of
        Section~\ref{sec:meth-c8-space}.
\end{enumerate}

\section{Continuation of this work}
\label{sec:concl-future}

The post-analysis of Campaign 9 shows that the formulation of the
problem and the surrogate-assisted loop can be refined in a new
campaign. Five changes to a new campaign have the most value, in decreasing order of their expected
effect on its results:

\begin{enumerate}
    \item \textbf{Eight-layer depletion in place of the assembly
        depletion:} the evaluator depletes the three-dimensional core in
        8 axial layers instead of the reflective assembly:
        \begin{itemize}
          \item \textbf{Reason:} the assembly depletion, with its
                axially uniform burnup, overestimates the cycle length
                of the core by 15 to 25 per cent, and no correction of
                it decides feasibility near the minimum cycle length
                (Section~\ref{sec:res-c9-axial}).
          \item \textbf{Setting:} $20\,000$ particles per batch over 90
                batches, 40 of them inactive, which gives a statistical
                error of about \SI{21}{d} on one cycle length, the
                measured \SI{15.1}{d} at 100 active batches scaled to
                50, below the \SI{29.1}{d} of the assembly depletion of
                Campaign 9, at about \SI{40}{min} per evaluation.
          \item \textbf{Expected effect:} the cycle-length constraint is
                evaluated on the burnup distribution the core develops,
                the surrogate is trained on that value, and the search
                places its front at the minimum cycle length of the core
                instead of 15 to 25 per cent above it.
          \item \textbf{Confirmation:} the front members are then
                depleted again at the setting of the post-analysis,
                $20\,000$ particles per batch over 160 batches, 60 of
                them inactive, as in the two-tier fidelity strategy of
                Section~\ref{sec:meth-fid-settings}, where the search
                uses a low-fidelity depletion, so the reduced fidelity
                selects the candidates and the higher fidelity confirms
                them.
        \end{itemize}
        \item \textbf{MTC boron limit as a constraint of the search:} the
        a posteriori re-execution of Section~\ref{sec:res-c9-acq} shows
        that, with the MTC boron limit as a constraint of the search, the
        batches of Campaign 9 contain no design above the limit and 5
        instead of 11 designs below \SI{1}{wt\%} of gadolinia, the region
        where the designs are infeasible on that limit and where the
        campaign spent 6 of its 36 infill evaluations and stage 2 of the
        continuation 10 of its 18. In a new campaign the limit of each
        lattice, or until it is measured the logarithmic law of
        Section~\ref{sec:res-c9-ceiling} inside its range, enters the
        constraints of the surrogate search.
        \item \textbf{Acquisition function for an optimum on a constraint:}
        the acquisition function is changed so that it still
        discriminates between candidates when the optimum is on a
        constraint, with one of two proposals:
        \begin{itemize}
          \item \textbf{Confidence bound inside the search:} the
                constraint $\mu_{ij} + \kappa\,\sigma_{ij} \le 0$ of safe
                exploration~\cite{sui_safeopt_2015} is applied inside the
                NSGA-II search, so every candidate the search returns
                already satisfies the margin, at the cost of a front kept
                about \SI{42}{d} above the minimum cycle length.
          \item \textbf{Probability of feasibility as a weight:} the
                uncertainty score is multiplied by the probability of
                feasibility of constrained Bayesian
                optimisation~\cite[Sec.~5.4.2]{forrester_book_2008},
                which keeps a design on the constraint at half weight
                instead of rejecting it. The re-execution shows that a
                ranking by the probability of feasibility alone is worse
                than the margin, so the probability weights the score
                and does not replace it.
        \end{itemize}
        \item \textbf{Operating conditions in the depletion:} a critical
        boron search at each burnup step in place of the constant
        \SI{1000}{ppm}, and thermal-hydraulic coupling to a subchannel
        solver in place of the fixed temperatures, which give the cycle
        length of an operated core in place of a lower bound, and the
        departure from nucleate boiling and fuel temperature margins of
        the candidates.
    \item \textbf{Shutdown margin as a constraint:} the cold zero-power
        core with every control rod assembly inserted except the one of
        highest worth, evaluated at beginning of life and at the
        operating maximum, in place of the controllability constraint
        at hot full power, which gives a licensing-style shutdown
        criterion.
\end{enumerate}
